\section{Step 5 Mix one octet and start eight bit steps}\label{crc-s05}

Load an unsigned octet, XOR it into the low byte of the raw state, and initialize the bit counter. These three operations prepare one byte; they do not yet advance the source. The first default byte is hexadecimal 31, so the mixed state is \texttt{ffffffce}.

\Needspace{10\baselineskip}
\subsection{Write and inspect}
Replace the marked temporary instructions in \texttt{b\_byte}. Retain surrounding labels and markers.

\MARSCodeLabel{MIPS assembly}
\begin{minted}[samepage=true]{mars-idea}
    lbu $t0, 0($a0)
    xor $v0, $v0, $t0
    li $t1, 8
\end{minted}

Use \texttt{s05-1}, \texttt{s05-2}. Prepare the exact memory inputs in a separate source file and pause at \texttt{kernel\_done}, before the caller stores the result.

\subsection{Expected observation}
The first check byte leaves \texttt{v0}=\texttt{ffffffce} and \texttt{t1}=8.

\Needspace{11\baselineskip}
\subsection{Check your understanding}
Why use an unsigned byte value here?

\begin{itemize}
  \item \textbf{A} It prevents source sign extension from changing higher state bits.
  \item \textbf{B} All files contain text.
  \item \textbf{C} It selects a larger buffer.
\end{itemize}
